\documentclass[11pt]{article}
\usepackage[margin=1in]{geometry}
\usepackage{amsmath,amssymb,amsthm}
\usepackage{newtxtext,newtxmath}
\usepackage{booktabs}
\usepackage{enumitem}
\usepackage{hyperref}
\usepackage{cleveref}

\newtheorem{definition}{Definition}
\newtheorem{proposition}{Proposition}
\newtheorem{theorem}{Theorem}
\newtheorem{corollary}{Corollary}
\newtheorem{remark}{Remark}

\title{Multi-Stakeholder Goodhart: A Study Guide\\[4pt]
\large From Intuition to Formalization}
\author{Working notes for RecSys paper}
\date{}

\begin{document}
\maketitle

\section{Level 1: The Intuition (Smart High Schooler)}

\textbf{Goodhart's Law:} ``When a measure becomes a target, it ceases to
be a good measure.''

Imagine a school improving education by tracking test scores. Teachers
start teaching to the test. Scores rise. Learning doesn't---and might
decline.

Three ingredients:
\begin{enumerate}
  \item You care about something hard to measure (real learning).
  \item You pick an easier proxy (test scores).
  \item You optimize hard for the proxy.
  \item The proxy diverges from what you actually care about.
\end{enumerate}

The proxy and the real thing overlap imperfectly. Gentle optimization
follows the overlap. Hard optimization finds and exploits the cracks.

\textbf{In recommendation systems:} A platform cares about user
wellbeing (hard to measure), so it optimizes engagement (clicks, likes).
Light optimization: relevant content gets engagement. Heavy
optimization: clickbait maximizes engagement but harms the user.

\textbf{The multi-stakeholder twist:} Now add a third party---society.
Society cares about diversity and reduced polarization. The platform
optimizing for one stakeholder (user engagement) may actively harm
another (societal welfare). The more precisely you optimize for one,
the worse it gets for the other. That's multi-stakeholder Goodhart.

%═══════════════════════════════════════════════════════════════
\section{Level 2: The Math (Undergrad)}
%═══════════════════════════════════════════════════════════════

\subsection{Single-Objective Goodhart}

Let $\pi$ be a policy (e.g., a recommendation strategy). Define:
\begin{align}
  U^*(\pi) &= \text{true utility (what we actually want)} \\
  \hat{U}(\pi) &= \text{proxy utility (what we optimize)}
\end{align}

We find $\pi^* = \arg\max_\pi \hat{U}(\pi)$ but care about $U^*(\pi^*)$.

\begin{definition}[Goodhart Effect]
A Goodhart effect occurs when $U^*(\pi^*) < U^*(\pi_0)$ for some
reference policy $\pi_0$---the optimized policy is \emph{worse} on the
true objective than the baseline.
\end{definition}

\textbf{When does this happen?} Write $\hat{U} = U^* + \epsilon$ where
$\epsilon$ captures misspecification. Think geometrically: $U^*$ and
$\hat{U}$ define two surfaces over policy space. Where they're aligned
(small $\epsilon$), optimizing $\hat{U}$ also climbs $U^*$. Where they
diverge, the optimizer finds a policy scoring high on $\hat{U}$ but
sitting in a valley of $U^*$.

\subsection{The Optimization Pressure Curve}

The critical variable is \textbf{optimization pressure} (how hard you
push toward $\hat{U}$). In our setting, this is controlled by $N$, the
number of preference pairs:

\begin{center}
\begin{tabular}{lll}
\toprule
Regime & $\hat{U}(\pi)$ & $U^*(\pi)$ \\
\midrule
$N$ small (underoptimized) & Low & Low \\
$N$ moderate (sweet spot) & Moderate & \textbf{Maximum} \\
$N$ large (overoptimized) & High & Declining \\
\bottomrule
\end{tabular}
\end{center}

At the sweet spot, the model has learned enough of the overlapping
signal to be useful, but not enough to precisely exploit the divergence.

\subsection{In Preference Learning}

A Bradley-Terry model learns weights $\mathbf{w} \in \mathbb{R}^D$ from
preference pairs. Each pair says ``content $a$ is preferred to content
$b$,'' and BT minimizes:
\begin{equation}
  \mathcal{L}_{\text{BT}} = -\log \sigma\!\left(
    \mathbf{w}^\top \phi(a) - \mathbf{w}^\top \phi(b)
  \right)
\end{equation}
where $\phi(x) \in \mathbb{R}^D$ is the feature vector (e.g., genre
vector for MovieLens) and $\sigma$ is the sigmoid.

If the preference labels come from a \emph{misspecified} utility
$\hat{\mathbf{w}}^\top \phi(x)$ instead of the true
$\mathbf{w}^{*\top} \phi(x)$:
\begin{itemize}
  \item BT with infinite data recovers $\hat{\mathbf{w}}$ exactly
    (Sun et al.'s order-consistency theorem).
  \item The learned policy selects content by
    $\arg\max_x \hat{\mathbf{w}}^\top \phi(x)$.
  \item The true utility of this selection is
    $\mathbf{w}^{*\top} \phi(\arg\max_x \hat{\mathbf{w}}^\top
    \phi(x))$.
\end{itemize}

More training pairs $\Rightarrow$ $\hat{\mathbf{w}}$ is learned more
precisely $\Rightarrow$ content selection is more aligned with the
\emph{wrong} direction $\Rightarrow$ true utility decreases.

%═══════════════════════════════════════════════════════════════
\section{Level 3: The Formal Framework (Grad Student)}
%═══════════════════════════════════════════════════════════════

\subsection{Skalse et al.'s Taxonomy (ICLR 2024)}

Let $R^*: \mathcal{S} \times \mathcal{A} \to \mathbb{R}$ be the true
reward, $\hat{R}$ a proxy. An agent trained to maximize $\hat{R}$
exhibits Goodhart when
$J^*(\pi_{\hat{R}}) < J^*(\pi_{R^*})$
where $J^*(\pi) = \mathbb{E}[\sum_t R^*(s_t, a_t) \mid \pi]$.

Four mechanisms:

\begin{enumerate}
  \item \textbf{Regressional:} $\hat{R}$ is an unbiased estimator of
    $R^*$ with noise. Optimizing $\hat{R}$ selects for positive noise
    (optimizer's curse). Effect is $O(\sigma^2/n)$.

  \item \textbf{Extremal:} $\hat{R}$ and $R^*$ agree over most of
    policy space but diverge in extreme regions. The optimizer is drawn
    to precisely these regions.

  \item \textbf{Causal:} $\hat{R}$ captures a correlate of $R^*$,
    not a cause. Optimization breaks the correlation.

  \item \textbf{Adversarial:} The proxy is deliberately constructed to
    exploit the optimizer. (Not relevant to our setting.)
\end{enumerate}

Skalse et al.\ prove: for any $\hat{R} \neq R^*$ and any MDP, there
exists a critical optimization threshold $\beta^*$ such that for
$\beta > \beta^*$, the optimized policy has lower true return than
the reference.

\subsection{Gao et al.'s Scaling Laws (ICML 2023)}

\textbf{The setup.} You can't cheaply get human preferences at scale, so
Gao et al.\ use a synthetic setup: a large ``gold-standard'' reward model
$R_{\text{gold}}$ plays the role of humans. A smaller ``proxy'' reward
model $R_{\text{proxy}}$ is trained on labels from $R_{\text{gold}}$.
Then a policy is optimized against $R_{\text{proxy}}$ using either
reinforcement learning (PPO) or best-of-$n$ sampling. The question: as
you optimize harder against the proxy, what happens to the gold score?

\textbf{The key equation.} For both RL and best-of-$n$, the gold RM
score as a function of KL divergence $d$ from the reference policy is
well-fit by:
\begin{equation}
  R_{\text{gold}}(d) = \alpha \sqrt{d} - \beta \cdot d
  \label{eq:gao}
\end{equation}
This has two terms with distinct interpretations:
\begin{itemize}
  \item $\alpha\sqrt{d}$: the \textbf{learning term}. Genuine alignment
    improvement. Grows as $\sqrt{d}$ --- diminishing returns from
    optimization. This corresponds to the model learning the shared
    signal between proxy and gold rewards.

  \item $\beta \cdot d$: the \textbf{overoptimization term}. Exploitation
    of proxy imperfection. Grows \emph{linearly} in $d$, so it
    eventually dominates the sublinear learning term.
\end{itemize}

The gold score peaks at:
\begin{equation}
  d^* = \frac{\alpha^2}{4\beta^2}
  \qquad\text{with peak gold score}\qquad
  R^*_{\text{gold}} = \frac{\alpha^2}{4\beta}
\end{equation}
Beyond $d^*$, every additional unit of optimization \emph{decreases}
the true reward. The optimizer is now in Goodhart territory.

\textbf{Why $\sqrt{d}$ vs.\ $d$?} The intuition: the beneficial
component of the proxy reward corresponds to the \emph{projection} of
the proxy onto the true reward direction. In a high-dimensional space,
this projection is a lower-dimensional signal that saturates quickly
($\sqrt{d}$). The harmful component---the orthogonal residual---has
more degrees of freedom and grows linearly. This is a geometric
argument: in $D$ dimensions, the projection onto a 1D truth direction
concentrates at rate $1/\sqrt{D}$, while the residual has $D-1$
dimensions to exploit.

\textbf{Scaling with reward model size.} As the proxy RM gets larger
(more parameters), $\alpha$ increases (better proxy $\Rightarrow$ more
genuine learning) and $\beta$ decreases (fewer exploitable flaws).
Both effects push $d^*$ rightward --- you can optimize harder before
hitting Goodhart. Gao et al.\ find smooth power-law scaling of
$\alpha$ and $\beta$ with RM parameter count.

\textbf{Scaling with dataset size.} More labels for training the proxy
RM also push $d^*$ rightward, but with diminishing returns. Doubling
the label dataset is less effective than doubling the RM size.

\textbf{RL vs.\ best-of-$n$.} Both follow Equation~\ref{eq:gao}, but
with different coefficients. The key finding: \emph{RL consumes much
more KL than best-of-$n$ for the same amount of optimization.} At the
same $d_{\text{KL}}$, best-of-$n$ achieves a higher gold score. RL
``wastes'' KL budget on distributional changes that don't improve
alignment.

\textbf{KL penalty $\approx$ early stopping.} Adding a KL penalty to
the RL objective does not change the $d$-vs-gold-score frontier. Any
point on the frontier can be reached by early stopping alone, without
the KL penalty. This is a surprising and somewhat discouraging result:
the standard RLHF regularization technique doesn't actually buy you
anything beyond what early stopping provides.

\textbf{Connection to our work.} Gao et al.\ parameterize
overoptimization by $d_{\text{KL}}$ (distance in policy space). We
parameterize by $N$ (number of preference pairs). The analogy:
\begin{center}
\begin{tabular}{ll}
\toprule
\textbf{Gao et al.} & \textbf{Our setting} \\
\midrule
KL divergence $d$ & Training pairs $N$ \\
Gold RM score & Hidden stakeholder utility \\
Proxy RM & Misspecified BT weights \\
Peak at $d^*$ & Sweet spot at $N^*$ \\
$\alpha\sqrt{d}$: learning & BT convergence toward true signal \\
$\beta \cdot d$: Goodhart & BT convergence toward misspec.\ direction \\
\bottomrule
\end{tabular}
\end{center}

Our $N^* \approx 50$ on MovieLens plays the same role as their $d^*$:
the point where learning saturates and overoptimization begins. A
natural question is whether our data admits a similar
$\alpha\sqrt{N} - \beta N$ functional form.

\subsection{Catastrophic Goodhart (Kwa et al., NeurIPS 2024)}

\textbf{The problem with KL regularization.} The standard defense
against reward hacking in RLHF is to penalize the KL divergence
between the optimized policy and the reference:
\begin{equation}
  \max_\pi \; \mathbb{E}_\pi[\hat{R}(x)] - \lambda \cdot
  D_{\text{KL}}(\pi \| \pi_{\text{ref}})
\end{equation}
The idea: by staying ``close'' to the reference policy (in KL sense),
you can't deviate too far from reasonable behavior, even if $\hat{R}$
is imperfect. Gao et al.'s work implicitly relies on this—their
functional form assumes KL controls the damage.

Kwa et al.\ show this assumption fails under a specific, realistic
condition.

\textbf{The key distinction: light-tailed vs.\ heavy-tailed error.}
Define the proxy error $\epsilon(x) = \hat{R}(x) - R^*(x)$. The
distribution of $\epsilon$ across responses $x$ drawn from the
reference policy can be:
\begin{itemize}
  \item \textbf{Light-tailed} (e.g., Gaussian): extreme errors are
    exponentially rare. $P(|\epsilon| > t) \leq e^{-ct^2}$.
  \item \textbf{Heavy-tailed} (e.g., Pareto, log-normal): extreme
    errors are polynomially rare. $P(|\epsilon| > t) \sim t^{-\alpha}$.
\end{itemize}

\textbf{Theorem (informal): Light tails $\Rightarrow$ KL works.}
When $\epsilon$ is light-tailed and independent of $R^*$, optimal
policies under decreasing KL penalties achieve arbitrarily high true
utility. The KL constraint successfully limits exploitation because
extreme proxy rewards (where $\epsilon$ is large) are exponentially
unlikely, so upweighting them costs exponential KL.

\textbf{Theorem (informal): Heavy tails $\Rightarrow$ KL fails
catastrophically.}
When $\epsilon$ is heavy-tailed, there exist policies that achieve:
\begin{equation}
  \mathbb{E}_\pi[\hat{R}] \to \infty
  \quad\text{while}\quad
  D_{\text{KL}}(\pi \| \pi_{\text{ref}}) \to 0
  \quad\text{and}\quad
  \mathbb{E}_\pi[R^*] \leq \mathbb{E}_{\pi_{\text{ref}}}[R^*]
\end{equation}
In words: the policy gets arbitrarily high proxy reward with
\emph{negligible} KL cost, while achieving \emph{no improvement}
in true utility. The KL penalty is useless.

\textbf{Why does this happen?} The key insight is about how KL
divergence interacts with tails. To achieve high expected proxy reward,
the policy needs to upweight responses where $\hat{R}$ is large. If
$\hat{R}$ is large because $\epsilon$ is large (not because $R^*$ is
large), this is Goodhart.

The KL cost of upweighting a response $x$ by factor $t$ is
$\sim \log t$. The proxy reward gain is $\sim \hat{R}(x)$. Under
heavy tails, there exist responses with $\hat{R}(x) \sim t^\alpha$
but probability $\sim t^{-\alpha}$ under the reference. Upweighting
by factor $t$ costs $\text{KL} \sim \log t$ but gains reward
$\sim t^\alpha$. As $t \to \infty$, the reward-to-KL ratio diverges.
The optimizer can always find more to exploit, and KL can never
constrain it.

Under light tails, the same responses have probability
$\sim e^{-ct^2}$, so upweighting by factor $t$ costs
$\text{KL} \sim ct^2$. The cost grows \emph{faster} than the reward,
so eventually KL wins.

\subsubsection*{Why KL Fails: The Geometry of Average vs.\ Concentrated}

Your intuition is correct: KL is an average measure, and the reward gain
is concentrated. Let's make this precise at three levels.

\textbf{Level 1: The arithmetic.} KL divergence is an \emph{expectation}
of the log-likelihood ratio:
\begin{equation}
  D_{\text{KL}}(\pi \| \pi_{\text{ref}}) = \sum_x \pi(x)
  \log\frac{\pi(x)}{\pi_{\text{ref}}(x)}
\end{equation}
Upweighting a response $x$ by factor $t$ (setting $\pi(x) = t \cdot
\pi_{\text{ref}}(x)$) contributes:
\begin{equation}
  \Delta \text{KL} = t \cdot \pi_{\text{ref}}(x) \cdot \log t
\end{equation}
The contribution to expected reward is:
\begin{equation}
  \Delta \text{Reward} = t \cdot \pi_{\text{ref}}(x) \cdot \hat{R}(x)
\end{equation}
The reward-to-KL ratio from this single response is
$\hat{R}(x) / \log t$. The penalty ($\log t$) is \textbf{concave} in
the upweighting factor. The reward ($\hat{R}(x)$) can be arbitrarily
large. So for any fixed KL budget, you can always find responses where
the reward gain exceeds the KL cost---\emph{if} such extreme-reward
responses exist with non-negligible probability.

That's the ``if'' that distinguishes light from heavy tails.

\textbf{Level 2: Why ``average'' fails against ``concentrated.''}
KL divergence measures the \emph{average} surprise when sampling from
$\pi$ but expecting $\pi_{\text{ref}}$. It's insensitive to what
happens at any \emph{specific} response---it only sees the aggregate.

The optimizer's reward gain is dominated by a \emph{few extreme
responses}---the ones where $\hat{R}(x)$ is large. Under heavy tails,
these responses have non-negligible probability. The optimizer
concentrates its probability budget on them.

The mismatch: KL spreads its penalty across the entire distribution
(every response that moved contributes). The reward concentrates its
gain on a few extreme points. An average penalty can't constrain a
concentrated gain---it's like trying to cap a spike with a blanket.

Under light tails, extreme responses are exponentially rare
($P(\hat{R} > t) \sim e^{-ct^2}$). Upweighting them enough to matter
for expected reward requires exponential probability inflation, which
costs exponential KL. The blanket is thick enough.

Under heavy tails, extreme responses are polynomially rare
($P(\hat{R} > t) \sim t^{-\alpha}$). Upweighting them costs logarithmic
KL but gains polynomial reward. The blanket is too thin.

\textbf{Level 3: The variational characterization.} KL divergence has a
precise optimality property (Donsker--Varadhan representation):
\begin{equation}
  D_{\text{KL}}(\pi \| \pi_{\text{ref}}) =
  \sup_f \left\{
    \mathbb{E}_\pi[f] - \log \mathbb{E}_{\pi_{\text{ref}}}[e^f]
  \right\}
\end{equation}
Rearranging: if $D_{\text{KL}}(\pi \| \pi_{\text{ref}}) \leq C$, then
for any function $f$:
\begin{equation}
  \mathbb{E}_\pi[f] \leq \log \mathbb{E}_{\pi_{\text{ref}}}[e^f] + C
  \label{eq:dv-bound}
\end{equation}
Apply this with $f = \hat{R}$ (the proxy reward):
\begin{equation}
  \mathbb{E}_\pi[\hat{R}]
  \leq \underbrace{\log \mathbb{E}_{\pi_{\text{ref}}}[e^{\hat{R}}]}
    _{\text{moment generating function}} + C
\end{equation}
When $\hat{R}$ is \textbf{light-tailed}, the MGF
$\mathbb{E}[e^{\hat{R}}]$ is finite, so the bound is meaningful:
KL budget $C$ limits the achievable reward.

When $\hat{R}$ is \textbf{heavy-tailed}, the MGF is \emph{infinite}:
$\mathbb{E}[e^{\hat{R}}] = \infty$. The bound becomes
$\mathbb{E}_\pi[\hat{R}] \leq \infty + C = \infty$. It says nothing.
KL has zero constraining power.

This is the deepest answer: KL divergence constrains expected values of
functions whose \textbf{moment generating functions are finite} under
the reference distribution. Heavy-tailed rewards violate this condition.
No finite KL budget can constrain them.

\textbf{Would other divergences help?}
\begin{itemize}
  \item \textbf{$\chi^2$ divergence:}
    $\chi^2(\pi \| \pi_{\text{ref}}) =
    \mathbb{E}_{\pi_{\text{ref}}}[(\pi/\pi_{\text{ref}} - 1)^2]$.
    Upweighting by $t$ costs $(t-1)^2$---\emph{quadratic}, not
    logarithmic. Much more effective against concentrated upweighting.
    But also much more restrictive (rules out many useful policies).

  \item \textbf{R\'enyi divergence of order $\alpha > 1$:}
    Interpolates between KL ($\alpha \to 1$) and max-divergence
    ($\alpha \to \infty$). Higher $\alpha$ penalizes tails more.
    Using R\'enyi with $\alpha > 1$ as the regularizer would provide
    stronger guarantees for heavy-tailed rewards, at the cost of more
    conservative policies.

  \item \textbf{Total variation:} $\text{TV}(\pi, \pi_{\text{ref}}) =
    \frac{1}{2} \|\pi - \pi_{\text{ref}}\|_1$. Bounds
    $|\mathbb{E}_\pi[f] - \mathbb{E}_{\pi_{\text{ref}}}[f]| \leq
    \|f\|_\infty \cdot \text{TV}$. Works for bounded $f$ but is
    vacuous for unbounded rewards.
\end{itemize}

The choice of divergence is a choice of \emph{which deviations you
penalize most}. KL penalizes average deviations (good for typical
behavior). $\chi^2$ penalizes concentrated deviations (good against
tail exploitation). The right choice depends on whether the threat is
diffuse or concentrated---which is exactly the light-tail / heavy-tail
distinction.

\textbf{Empirical finding: current reward models appear light-tailed.}
Kwa et al.\ test real reward models (Starling 7B, Pythia 1.4B) and find
error distributions that are ``approximately normal and, therefore,
light-tailed.'' Current RLHF systems are probably safe from
catastrophic Goodhart---but future systems with more powerful reward
models and optimization might not be.

\textbf{Connection to our work.} Our setting has a specific structure
that maps onto this framework:
\begin{itemize}
  \item Our ``proxy'' is the misspecified user utility
    $\hat{\mathbf{w}}_u^\top \phi(x)$.
  \item Our ``true reward'' is the diversity utility
    $\mathbf{w}_d^\top \phi(x)$.
  \item The ``error'' is $\epsilon(x) = (\hat{\mathbf{w}}_u -
    \mathbf{w}_d)^\top \phi(x)$, which is a \emph{linear function}
    of fixed genre features.
  \item Since genre features are bounded binary vectors, $\epsilon$
    is \emph{light-tailed} by construction. Catastrophic Goodhart
    does not apply.
\end{itemize}

However, our Goodhart effect is \emph{not} mitigated by this. Why?
Because we don't use KL regularization. Our optimizer (BT preference
learning) has no KL penalty---it simply converges toward the direction
specified by the training labels. The relevant parameter is $N$ (data
quantity), not KL divergence. Kwa et al.'s result is about the limits
of KL as a \emph{defense}; our result is about the limits of data as an
\emph{offense}---more data with wrong specification amplifies the error
regardless of distributional tail behavior.

Let's unpack this distinction and its implications.

\subsubsection*{Defense vs.\ Offense: Two Different Failure Modes}

\textbf{Kwa et al.'s world (defense failure):} There is a fixed,
imperfect proxy reward model $\hat{R}$. A policy optimizer (PPO) tries
to maximize $\hat{R}$ subject to a KL constraint. The question is
whether the constraint (the defense) is strong enough to prevent the
optimizer from reaching regions where $\hat{R} \neq R^*$. Under heavy
tails: no. Under light tails: yes, but the protected region scales
with the KL budget.

The causal chain is: optimizer pushes policy $\to$ KL tries to stop it
$\to$ heavy tails let the optimizer slip through.

\textbf{Our world (offense failure):} There is no policy optimizer. BT
preference learning is a \emph{supervised} learner---it fits weights to
preference labels via gradient descent. It has no ``pressure'' to push
beyond the training signal; it simply converges toward the direction
implied by the labels. The ``offense'' is the \emph{data itself}: each
preference pair is a gradient step toward the misspecified utility
direction.

The causal chain is: more data $\to$ BT converges more precisely $\to$
learned weights align with misspecified direction $\to$ content
selection changes $\to$ hidden stakeholder degrades.

There is nothing to ``defend against'' because the attack isn't
optimization pressure---it's specification error baked into the training
signal. KL regularization penalizes how far the policy moves; but in our
setting the policy moves \emph{because the data tells it to}, and more
data tells it more precisely.

\subsubsection*{What If We Switched to Softmax Sampling?}

Suppose we replaced deterministic top-$K$ selection with probabilistic
softmax sampling:
\begin{equation}
  P(x \mid \mathbf{w}) = \frac{\exp(\mathbf{w}^\top \phi(x) / \tau)}
    {\sum_{x'} \exp(\mathbf{w}^\top \phi(x') / \tau)}
\end{equation}
Now we have a stochastic policy and can define:
\begin{equation}
  D_{\text{KL}}(P(\cdot \mid \hat{\mathbf{w}}) \| P(\cdot \mid
  \mathbf{w}^*))
\end{equation}
Could we then use KL to prevent Goodhart?

\textbf{First observation:} We're in the light-tailed regime.
Our content features $\phi(x)$ are bounded binary genre vectors, so
$\mathbf{w}^\top \phi(x)$ is bounded for any $\mathbf{w}$. The softmax
distribution has bounded log-ratios. By Kwa et al.'s Theorem~4 (light
tails), KL \emph{would} provide a meaningful constraint. No catastrophic
Goodhart.

\textbf{Second observation:} KL would constrain but not eliminate the
effect. Adding KL between the learned and reference softmax policies is
equivalent to limiting how far the BT-learned weights can stray from
some reference weights. But this is just \emph{early stopping} by
another name---exactly Gao et al.'s finding that KL penalty
$\approx$ early stopping.

In our $N$-parameterization, this means: KL regularization would
formalize the sweet spot at $N^*$. It would let you control \emph{where}
on the learning-vs-Goodhart curve you operate. But the curve still
exists---the tension between user and diversity utility is structural,
not an artifact of the optimization method.

Concretely: with softmax + KL, you'd trade the parameter $N$ (data
quantity) for the parameter $\lambda$ (KL penalty strength). At high
$\lambda$ (strong penalty): the policy stays near the reference,
learning little---equivalent to low $N$. At low $\lambda$ (weak
penalty): the policy converges to the misspecified direction---equivalent
to high $N$. The Goodhart curve is reparameterized, not removed.

\textbf{Third observation:} Where would the reference weights come from?
In RLHF, the reference is the pretrained model (before fine-tuning).
In our setting, the reference would need to be ``correct'' stakeholder
weights---but if we had correct weights, we wouldn't need BT training at
all. Using random/initial weights as reference penalizes ALL learning
(good and bad), which is wasteful. The fundamental problem is that KL
penalizes \emph{distance} from reference, not \emph{direction} of
movement. If the movement direction is wrong (misspecified utility),
KL can slow you down but can't redirect you.

\subsubsection*{Heavy Tails, Extremal Goodhart, and Our Setting}

The four Goodhart mechanisms and the tail behavior form a
two-dimensional picture. Let's connect them:

\textbf{Regressional Goodhart} doesn't care about tails. It's about
noise in the proxy estimate, and it gets \emph{better} with more data
(the noise averages out). Tail behavior is irrelevant because the effect
is at the \emph{center} of the distribution, not the tails.

\textbf{Extremal Goodhart} is about what happens when the optimizer
reaches regions where proxy and truth diverge. Two questions:
\begin{enumerate}
  \item \textbf{Do such regions exist?} Yes, whenever proxy $\neq$
    truth. In our setting: the region where content is user-optimal but
    diversity-poor exists because $\cos(\mathbf{w}_u, \mathbf{w}_d) < 0$.
    These are movies in popular genres with high ratings---they exist in
    any realistic content pool.

  \item \textbf{Does the optimizer reach them?} This is where tails
    matter---but \emph{only for policy optimizers with KL constraints.}
    Under heavy tails, KL can't stop the optimizer from reaching the
    extremes. Under light tails, KL bounds how far the optimizer gets.
\end{enumerate}

\textbf{In our setting, tail behavior is irrelevant to the extremal
mechanism.} Here's why:

Our ``optimizer'' is BT preference learning, which converges to the
direction specified by the training labels. It doesn't search over a
policy space looking for extreme responses. It learns a weight vector
that ranks content. The content pool is \emph{finite} (3{,}347 movies
on ML-1M). The ``extreme'' items---movies that are user-optimal but
diversity-poor---are just specific movies in the pool. They're not in
any distributional tail; they're fixed items.

As $N$ increases, BT doesn't discover new extremes. It learns to
\emph{rank the existing items more accurately} according to the
misspecified utility. At $N = 25$, the ranking is noisy and the
top-$K$ selection includes a mix of items. At $N = 2000$, the ranking
precisely identifies the user-optimal items and puts them at the top.

The mechanism is:
\begin{equation*}
  \text{More data} \to \text{more precise ranking} \to
  \text{top-}K \text{ concentrates on misspec-optimal items} \to
  \text{diversity degrades}
\end{equation*}

This is extremal Goodhart---the optimizer reaches the extreme region
where proxy and truth diverge. But it reaches it through
\emph{convergence of a supervised learner}, not through
\emph{exploitation of distributional tails}. The finite content pool
removes any role for tail behavior.

\textbf{To summarize the relationship:}
\begin{center}
\begin{tabular}{p{4cm}p{4cm}p{4cm}}
\toprule
& \textbf{RLHF setting} & \textbf{Our setting} \\
\midrule
\textbf{What causes Goodhart?}
  & Policy optimizer finds extreme responses where $\hat{R} \neq R^*$
  & BT converges to misspecified direction; top-$K$ concentrates on
    wrong content \\[6pt]
\textbf{Mechanism}
  & Extremal: optimizer reaches tail of distribution
  & Extremal: supervised learner precisely learns wrong ranking
    over finite pool \\[6pt]
\textbf{Role of tail behavior}
  & Heavy tails $\Rightarrow$ KL can't prevent reaching extremes
  & Irrelevant: content pool is finite, no distributional tails \\[6pt]
\textbf{What controls severity?}
  & KL budget / tail heaviness
  & $N$ (data quantity) / $\cos(\mathbf{w}_{\text{target}},
    \mathbf{w}_{\text{hidden}})$ \\[6pt]
\textbf{Defense}
  & KL penalty (limited by tails)
  & Early stopping on $N$ / independent selection dim ($\delta$) \\
\bottomrule
\end{tabular}
\end{center}

Both settings exhibit extremal Goodhart. The \emph{channel} through
which the optimizer reaches the extreme region is different: policy
optimization in RLHF, supervised convergence in BT preference learning.
The tail behavior of the reward distribution governs the first channel
but is irrelevant to the second. In both cases, the core issue is the
same: the optimizer finds exactly what it's told to find, and if what
it's told is wrong, precision makes things worse.

\subsection{What None of These Address}

\begin{itemize}
  \item \textbf{Multiple objectives with known geometry.} All above
    treat one proxy vs.\ one true reward. None study settings where
    multiple stakeholders have utility vectors with known cosine
    similarity.

  \item \textbf{Data-dependent thresholds.} Gao et al.\ parameterize
    by KL divergence, not by the number of preference pairs $N$. The
    sweet spot in $N$ (our finding: $N \approx 50$ on MovieLens) is not
    studied.

  \item \textbf{Selection mechanism independence.} Our rank stability
    finding ($= 1.0$ for individual parameter perturbation) initially
    suggested the Pareto frontier ``buffers'' against Goodhart. The
    actual mechanism is more specific---see \cref{sec:pareto-buffer}
    below.

  \item \textbf{Low-rank weight space.} Our Phase~9 finding that
    the weight space is effectively 3-dimensional (projection cosine
    $> 0.97$) constrains the Goodhart dynamics---the
    misspecification can only push within a low-dimensional subspace.
\end{itemize}

\subsection{Why KL Divergence Doesn't Appear in Our Formulation}
\label{sec:why-no-kl}

Every paper above parameterizes overoptimization by KL divergence
between policies. We parameterize by $N$ (preference pairs). This isn't
an oversight---it reflects a structural difference in our optimization
pipeline.

\textbf{In RLHF:} The policy $\pi(y \mid x)$ is a distribution over
responses. KL divergence $D_{\text{KL}}(\pi \| \pi_{\text{ref}})$
measures how far the optimized distribution has moved from the
reference. The optimizer (PPO) is explicitly regularized by KL.

\textbf{In our setting:} The ``policy'' is a deterministic function:
given weight vector $\mathbf{w}$, select top-$K$ content by
$\mathbf{w}^\top \phi(x)$. There is no stochastic policy, no reference
distribution, and no KL penalty. The BT optimizer simply minimizes
preference loss on $N$ pairs until convergence.

\textbf{What plays the role of KL?} In our setting, the \textbf{cosine
distance} between the learned and true weight vectors:
\begin{equation}
  d_{\cos} = 1 - \cos(\hat{\mathbf{w}}(N), \mathbf{w}^*)
\end{equation}
serves the same conceptual role as $d_{\text{KL}}$. Both measure ``how
far the optimizer has moved from the reference.'' The difference:
\begin{itemize}
  \item $d_{\text{KL}}$ measures distance in \emph{policy distribution}
    space (infinite-dimensional, information-geometric).
  \item $d_{\cos}$ measures distance in \emph{weight direction} space
    ($D$-dimensional, Euclidean geometry on the unit sphere).
\end{itemize}

\textbf{Why this matters for Goodhart.} In RLHF, KL controls
optimization pressure---larger KL budget $\Rightarrow$ more
optimization $\Rightarrow$ more Goodhart. In our setting, $N$ controls
optimization pressure---more pairs $\Rightarrow$ BT converges more
precisely $\Rightarrow$ $d_{\cos}$ between $\hat{\mathbf{w}}(N)$ and
the \emph{misspecified} target increases (the model moves further from
the true direction toward the wrong direction).

The Gao et al.\ functional form $R_{\text{gold}}(d) = \alpha\sqrt{d} -
\beta d$ could potentially be re-expressed in our setting as
$U_{\text{hidden}}(d_{\cos}(N))$, where $d_{\cos}(N)$ captures the
BT convergence trajectory. The intermediate variable $d_{\cos}$ would
bridge between $N$ (our control parameter) and
$U_{\text{hidden}}$ (our outcome), with the $\sqrt{\cdot}$ term
capturing the learning phase and the linear term capturing the
Goodhart phase.

\textbf{Could we add KL to our setting?} Yes---by replacing hard
top-$K$ selection with softmax sampling:
\begin{equation}
  P(x \mid \mathbf{w}) = \frac{\exp(\mathbf{w}^\top \phi(x) / \tau)}
    {\sum_{x'} \exp(\mathbf{w}^\top \phi(x') / \tau)}
\end{equation}
Then $D_{\text{KL}}(P(\cdot \mid \hat{\mathbf{w}}) \| P(\cdot \mid
\mathbf{w}^*))$ is well-defined and connects directly to Gao/Kwa's
framework. For small weight perturbations, this KL relates to cosine
distance via the Fisher information of the softmax. This is a potential
direction for formalization---but it changes the content selection
mechanism from deterministic to stochastic, which may affect the
Goodhart dynamics.

\subsection{The ``Pareto Buffer'': What It Actually Is}
\label{sec:pareto-buffer}

Our rank stability $= 1.0$ finding initially suggested the Pareto
frontier ``buffers'' against Goodhart. On closer examination, the
mechanism is more specific and less about Pareto geometry than about
\textbf{selection mechanism independence}.

\textbf{Two ways to select an operating point:}

\begin{enumerate}
  \item \textbf{Selection via misspecified weights (Goodhart possible).}
    Use learned weights as the content scorer. The misspecification
    directly determines what gets recommended. More precise
    misspecification $\Rightarrow$ more precisely wrong selection
    $\Rightarrow$ hidden stakeholder degrades. This is our Phase~6
    Goodhart experiment.

  \item \textbf{Selection via independent parameter (buffer exists).}
    Pick the operating point using the diversity weight $\delta$, which
    is \emph{independent} of any stakeholder's utility weights.
    Perturbing weights shifts the utility \emph{values} at each $\delta$
    but doesn't change \emph{which} $\delta$ you choose. This is the
    rank stability experiment.
\end{enumerate}

\textbf{Why rank stability $= 1.0$:} The diversity weight $\delta$
controls the tradeoff between engagement and topic diversity. It's an
independent control knob---not derived from any stakeholder utility. When
you perturb a single utility weight, the utility \emph{values} at each
$\delta$ shift, but the dominance relations among $\delta$-points don't
change. The same 7 of 21 $\delta$-points remain Pareto-optimal.

\textbf{When the buffer fails:}
\begin{itemize}
  \item \textbf{Correlated perturbation:} Simultaneously perturbing
    multiple weights changes utility values by different amounts at
    different $\delta$-points, breaking the dominance ordering. Rank
    stability drops below 1.0.
  \item \textbf{Misspecified scorer:} If the learned (misspecified)
    weights are used as the content \emph{scorer} (not just the
    evaluator), the selection itself changes. $\delta$ still controls
    the diversity-engagement tradeoff, but the engagement component
    now points in the wrong direction.
\end{itemize}

\textbf{Connection to Gao et al.'s KL $\approx$ early stopping
result:} KL regularization in RLHF is supposed to be an independent
control (constraining how far the policy moves), but Gao et al.\ found
it's redundant with the optimization axis---KL penalty doesn't change
the KL-vs-gold-score frontier. Our diversity weight $\delta$ succeeds
where KL fails because $\delta$ is \emph{genuinely orthogonal} to the
stakeholder utility weights. It controls content \emph{diversity} (a
structural property of the recommendation set), not content
\emph{quality} (which depends on the utility direction). The Pareto
structure helps not because it's ``Pareto'' but because it provides a
control dimension that the optimizer can't corrupt.

\textbf{The real insight:} Separating the selection mechanism
(diversity weight) from the specification mechanism (stakeholder
utility weights) provides protection against individual
misspecification. This is a design principle, not a mathematical
theorem about Pareto frontiers. It suggests that recommendation
systems should maintain at least one control dimension (like diversity)
that is defined structurally (e.g., set-level entropy) rather than
derived from stakeholder utility weights.

%═══════════════════════════════════════════════════════════════
\section{Our Empirical Results}
%═══════════════════════════════════════════════════════════════

\subsection{The Multi-Stakeholder Goodhart Curve}

Three stakeholders: user ($\mathbf{w}_u$), platform ($\mathbf{w}_p$),
diversity ($\mathbf{w}_d$), with:
\begin{equation}
  \cos(\mathbf{w}_u, \mathbf{w}_d) \approx -0.3, \quad
  \cos(\mathbf{w}_u, \mathbf{w}_p) \approx 0.95
\end{equation}

We perturb user weights: $\hat{\mathbf{w}}_u = \mathbf{w}_u \odot
(1 + \mathcal{N}(0, \sigma))$ with $\sigma = 0.3$, train BT on
$N$ preference pairs labeled by $\hat{\mathbf{w}}_u$, and measure
true diversity utility of the selected content.

\begin{center}
\begin{tabular}{rcc}
\toprule
$N$ & User utility (target) & Diversity utility (hidden) \\
\midrule
25 & 18.68 & $-6.70$ \\
50 & 19.44 & $\mathbf{-6.34}$ \textit{(peak)} \\
100 & 20.08 & $-7.24$ \\
200 & 20.81 & $-7.63$ \\
500 & 21.24 & $-8.38$ \\
2000 & 21.33 & $-8.98$ \\
\bottomrule
\end{tabular}
\end{center}

User utility increases monotonically ($+14\%$). Diversity utility peaks
at $N=50$, then degrades $42\%$ by $N=2000$. This is textbook
\textbf{extremal Goodhart}: the misspecified utility and the hidden
utility are anti-correlated, so more precise optimization of one
systematically harms the other.

\subsection{Key Empirical Findings}

\begin{enumerate}
  \item \textbf{The sweet spot exists and is data-dependent.}
    $N \approx 50$ on MovieLens (genre features, 19-dim).
    $N \approx 100$ on synthetic (action features, 18-dim).
    The sweet spot is where BT has learned the overlapping signal
    but hasn't yet converged to the misspecified direction.

  \item \textbf{The mechanism is utility transfer, not noise.}
    As $N$ increases, the content selection shifts from
    ``approximately random'' to ``precisely user-optimal,'' which is
    anti-optimal for diversity (cos $< 0$). This is systematic, not
    regressional.

  \item \textbf{Strategy 1 (better specification) confirms the
    mechanism.} Even with $\sigma = 0$ (correct specification), the
    user-optimal scorer selects content that harms diversity. Goodhart
    amplifies a tension that exists even under correct specification.

  \item \textbf{The weight space is low-rank.} Phase 9 showed projection
    cosine $> 0.97$ for scalarized $\to$ per-stakeholder span. The
    Goodhart dynamics are confined to a 3D subspace of the 19D genre
    feature space.
\end{enumerate}

%═══════════════════════════════════════════════════════════════
\section{Literature Map}
%═══════════════════════════════════════════════════════════════

\begin{center}
\begin{tabular}{lp{6cm}l}
\toprule
\textbf{Question} & \textbf{Status} & \textbf{Reference} \\
\midrule
Does Goodhart occur in RL? & Answered (yes, provably)
  & Skalse et al.\ 2024 \\
Scaling laws for overoptimization? & Answered (KL-based)
  & Gao et al.\ 2023 \\
Does KL regularization help? & Answered (no, heavy tails)
  & Kwa et al.\ 2024 \\
BT preference consistency? & Answered (order-consistent)
  & Sun et al.\ 2025 \\
Multi-stakeholder Pareto? & Framework exists
  & Burke \& Abdollahpouri 2017 \\
\midrule
Multi-stakeholder Goodhart? & \textbf{Open --- our work} & \\
$N$ as Goodhart modulator? & \textbf{Open --- our work} & \\
Stakeholder opposition ($\cos < 0$)? & \textbf{Open --- our work} & \\
Pareto as Goodhart buffer? & \textbf{Open --- our work} & \\
Low-rank weight constraints? & \textbf{Open --- our work} & \\
\bottomrule
\end{tabular}
\end{center}

%═══════════════════════════════════════════════════════════════
\section{Toward Formalization}
%═══════════════════════════════════════════════════════════════

\subsection{Setup}

Let $K$ stakeholders have utility directions
$\mathbf{w}_1, \ldots, \mathbf{w}_K \in \mathbb{R}^D$. A content pool
$\{\phi(x)\}_{x=1}^M$ with feature vectors $\phi(x) \in \mathbb{R}^D$.
A BT model learns $\hat{\mathbf{w}}(N)$ from $N$ preference pairs
labeled by a (possibly misspecified) utility.

Content selection: $S(\mathbf{w}) = \text{top-}k$ by
$\mathbf{w}^\top \phi(x)$.

Hidden stakeholder utility:
$U_{\text{hidden}}(S) = \frac{1}{k} \sum_{x \in S} \mathbf{w}_{\text{hidden}}^\top \phi(x)$.

\subsection{The Goodhart Condition}

\begin{proposition}[Informal]
More data hurts the hidden stakeholder when the learned direction
moves \emph{away} from the hidden utility direction:
\begin{equation}
  \frac{\partial U_{\text{hidden}}(S(\hat{\mathbf{w}}(N)))}{\partial N} < 0
  \quad \Longleftrightarrow \quad
  \nabla_{\mathbf{w}} U_{\text{hidden}} \cdot
  \frac{\partial \hat{\mathbf{w}}}{\partial N} < 0
\end{equation}
The left side is the chain rule: how hidden utility changes with $N$,
mediated by how the learned weights change the selection and how the
selection affects hidden utility.
\end{proposition}

\textbf{Unpacking the two terms:}

\begin{itemize}
  \item $\nabla_{\mathbf{w}} U_{\text{hidden}}$: How does hidden utility
    change as the scorer direction shifts? This is approximately
    $\mathbf{w}_{\text{hidden}}$ projected onto the content feature
    space --- it points toward content that's good for the hidden
    stakeholder.

  \item $\frac{\partial \hat{\mathbf{w}}}{\partial N}$: The direction
    BT converges toward as $N$ increases. With misspecified labels, this
    points toward $\hat{\mathbf{w}}_{\text{misspec}}$ --- the
    misspecified utility direction.
\end{itemize}

\textbf{The Goodhart condition:}
\begin{equation}
  \cos(\mathbf{w}_{\text{hidden}}, \hat{\mathbf{w}}_{\text{misspec}}) < 0
\end{equation}
More data hurts when the hidden stakeholder's utility direction is
anti-correlated with the direction BT is converging toward. In our
setting: $\cos(\mathbf{w}_{\text{diversity}},
\hat{\mathbf{w}}_{\text{user}}) \approx -0.3 < 0$. Goodhart is
predicted, and confirmed empirically.

\subsection{Numerical Validation (Phase 10)}

The direction condition is validated on 4 data points across 2 datasets,
with zero violations:

\begin{center}
\begin{tabular}{llrcl}
\toprule
\textbf{Dataset} & \textbf{Hidden} & $\cos(\mathbf{w}_u, \mathbf{w}_h)$
  & \textbf{Predicted} & \textbf{Actual} \\
\midrule
ML-100K & platform & $+0.956$ & Improve & $+14\%$ \checkmark \\
ML-100K & diversity & $-0.313$ & Degrade & $-42\%$ \checkmark \\
ML-1M & platform & $+0.926$ & Improve & $+14\%$ \checkmark \\
ML-1M & diversity & $-0.155$ & Degrade & $-46\%$ \checkmark \\
\bottomrule
\end{tabular}
\end{center}

\textbf{Direction condition holds perfectly.} Every stakeholder with
$\cos < 0$ degrades with $N$; every stakeholder with $\cos > 0$
improves. Platform ($\cos \approx +0.95$) serves as the positive
control: it improves monotonically because BT converges toward a
direction \emph{aligned} with platform utility.

\textbf{Magnitude condition is more nuanced.} ML-1M shows more
degradation (46\%) than ML-100K (42\%) despite \emph{less}
anti-correlation ($-0.155$ vs $-0.313$). The likely explanation:
ML-1M has more content (3{,}347 vs 1{,}305 movies), giving the
optimizer more user-optimal/diversity-poor items to concentrate on.
Degradation severity depends on both the angle between utility
directions \emph{and} the geometry of the content pool---the angle
determines the \emph{direction} of degradation, and the content pool
determines how much utility can be transferred in that direction.

\textbf{Statement of the validated condition:}

\begin{proposition}[Multi-Stakeholder Goodhart Direction Condition]
\label{prop:direction}
In a linear BT preference learning setting with $K$ stakeholders and a
finite content pool, let $\hat{\mathbf{w}}_t$ be the (possibly
misspecified) target utility learned from $N$ preference pairs, and
$\mathbf{w}_h$ be a hidden stakeholder's true utility. Then:
\begin{itemize}
  \item If $\cos(\hat{\mathbf{w}}_t, \mathbf{w}_h) < 0$:
    $U_h(S(\hat{\mathbf{w}}_t(N)))$ is eventually decreasing in $N$.
  \item If $\cos(\hat{\mathbf{w}}_t, \mathbf{w}_h) > 0$:
    $U_h(S(\hat{\mathbf{w}}_t(N)))$ is eventually increasing in $N$.
\end{itemize}
Validated on 3 datasets (synthetic, ML-100K, ML-1M) with 0 violations.
\end{proposition}

\subsection{Open Questions for Formalization}

\begin{enumerate}
  \item \textbf{Rate of degradation.} How fast does
    $U_{\text{hidden}}$ decline with $N$? Is it
    $O(1/\sqrt{N})$ (BT convergence rate) or slower/faster?

  \item \textbf{Sweet spot location.} Can we predict $N^*$ (the
    peak of $U_{\text{hidden}}$) from the cosine similarity and
    feature dimensionality? Our data: $N^* \approx 50$ on ML-100K
    (19-dim), $N^* \approx 25$ on ML-1M (19-dim, more content).

  \item \textbf{Pareto buffer.} Why does rank stability $= 1.0$ for
    individual parameter perturbation but $< 1.0$ for correlated
    perturbation? Can we characterize the perturbation directions
    that the Pareto frontier absorbs vs.\ amplifies?

  \item \textbf{Low-rank constraint.} With $K$ stakeholders in $D$
    dimensions, the effective weight space is $K$-dimensional
    (our Phase~9 finding). Does this bound the Goodhart severity?
    Intuitively, a 3D Goodhart in 19D is less severe than a 19D
    Goodhart --- fewer directions to exploit.

  \item \textbf{The ``fix spec first'' prescription.} Our data shows:
    data helps when specification is correct, data hurts when wrong.
    Can we formalize the \emph{crossover}: at what $\sigma$
    (misspecification level) does additional data switch from
    beneficial to harmful?
\end{enumerate}

%═══════════════════════════════════════════════════════════════
\section{References}
%═══════════════════════════════════════════════════════════════

\begin{enumerate}[label={[\arabic*]}]
  \item Skalse, Karwowski, Hayman, Bai, Kiendlhofer, Griffin.
    \textit{Goodhart's Law in Reinforcement Learning.} ICLR 2024.

  \item Gao, Schulman, Hilton.
    \textit{Scaling Laws for Reward Model Overoptimization.} ICML 2023.

  \item Kwa, Cunningham, Hoang.
    \textit{Catastrophic Goodhart: Regularizing RLHF with KL Divergence
    Does Not Mitigate Heavy-Tailed Reward Misspecification.} NeurIPS 2024.

  \item Sun, Shen, Ton.
    \textit{Rethinking Bradley-Terry Models in Preference-Based Reward
    Modeling.} ICLR 2025.

  \item Burke, Abdollahpouri.
    \textit{Towards Multi-Stakeholder Utility Evaluation of Recommender
    Systems.} VAMS Workshop 2017.

  \item Goodhart, C.A.E.
    \textit{Problems of Monetary Management: The U.K.\ Experience.}
    Papers in Monetary Economics, Reserve Bank of Australia, 1975.
\end{enumerate}

%═══════════════════════════════════════════════════════════════
\newpage
\section{Appendix: Unpacking Skalse et al.'s Four Mechanisms}
%═══════════════════════════════════════════════════════════════

Skalse et al.\ (ICLR 2024) identify four ways that optimizing a proxy
reward $\hat{R}$ can decrease the true reward $R^*$. These aren't
competing theories---they're four \emph{mechanisms} that can co-occur.
Understanding each one requires thinking about WHY the proxy and reality
diverge, not just THAT they diverge.

\subsection{Regressional Goodhart}

\textbf{The everyday version.} You're hiring and use SAT scores as a
proxy for job performance. Among 1000 applicants, the person with the
highest SAT score is more likely to have gotten lucky on test day than
to be genuinely the smartest. By selecting the extreme, you're selecting
for noise.

\textbf{The math.} Suppose the proxy is a noisy measurement of truth:
\begin{equation}
  \hat{R}(\pi) = R^*(\pi) + \epsilon(\pi), \quad
  \epsilon \sim \mathcal{N}(0, \sigma^2)
\end{equation}
where $\epsilon$ is independent noise that varies across policies.
The optimizer selects:
\begin{equation}
  \pi^* = \arg\max_\pi \left[ R^*(\pi) + \epsilon(\pi) \right]
\end{equation}
This selects for policies where both $R^*$ and $\epsilon$ happen to be
large. The expected true reward of the selected policy is:
\begin{equation}
  \mathbb{E}[R^*(\pi^*)] = \max_\pi R^*(\pi) - \underbrace{
    \mathbb{E}\!\left[\epsilon(\pi^*) \mid \pi^* = \arg\max(\hat{R})\right]
  }_{\text{optimizer's curse} > 0}
\end{equation}
The stronger the optimization (searching over more policies, or
optimizing more aggressively), the more you select for noise. This is
exactly the \textbf{winner's curse} from auction theory.

\textbf{Key properties:}
\begin{itemize}
  \item The proxy is \emph{unbiased}---$\mathbb{E}[\hat{R}] = R^*$ for
    any fixed policy. The problem is only in the \emph{selection}.
  \item The effect scales as $O(\sigma^2 / n)$ where $n$ is the number
    of samples used to estimate $\hat{R}$. More data reduces it.
  \item This is the \emph{mildest} form of Goodhart. It vanishes with
    perfect measurement.
\end{itemize}

\textbf{In our setting:} Regressional Goodhart would occur if BT
training noise caused us to select ``lucky'' content that happened to
score high due to noise in the learned weights. With 18--19 features and
2000+ preference pairs, this effect is small (our held-out accuracy is
within 2\% of training accuracy). It's not the dominant mechanism.

\subsection{Extremal Goodhart}

\textbf{The everyday version.} A doctor uses blood pressure as a proxy
for heart health. For most patients, lower blood pressure = healthier
heart. But at the \emph{extreme}---a patient with blood pressure of
60/40---the correlation breaks. The proxy and reality agreed in the
normal range but diverge at the tail.

\textbf{The math.} The proxy and truth are correlated but not identical:
\begin{equation}
  \hat{R}(\pi) = f(R^*(\pi)) + g(\pi)
\end{equation}
where $f$ is monotonically increasing (the correlation) and $g$ captures
features of $\pi$ that affect the proxy but not the truth. In the normal
range of $\pi$, $f$ dominates. At the extremes, $g$ can dominate.

The optimizer pushes $\pi$ into extreme regions where $g(\pi)$ is large.
The proxy says this is a great policy (high $\hat{R}$). The truth says
otherwise.

\begin{remark}
This differs from regressional Goodhart in a crucial way: the divergence
is \emph{systematic}, not random. The proxy is biased in extreme regions,
not just noisy. More data doesn't fix it---the bias is structural.
\end{remark}

\textbf{Key properties:}
\begin{itemize}
  \item Requires the optimizer to reach extreme regions of policy space.
    Weak optimization may never trigger it.
  \item The divergence is \emph{predictable} if you know how $g$
    behaves at extremes.
  \item More data makes it \emph{worse} (the optimizer more precisely
    finds the extreme region). This is the ``U-shaped curve'' in true
    performance.
\end{itemize}

\textbf{In our setting:} This is the \textbf{dominant mechanism}. The
``proxy'' is the misspecified user utility $\hat{\mathbf{w}}_u$. The
``truth'' is the diversity utility $\mathbf{w}_d$. They're correlated
(both are functions of genre features) but anti-correlated at the
extremes ($\cos \approx -0.3$). As BT learns $\hat{\mathbf{w}}_u$ more
precisely (more data), it pushes content selection into regions that are
\emph{great} for the misspecified user but \emph{terrible} for
diversity. The divergence is systematic: it's driven by the angle
between utility directions, not by noise.

\textbf{Why this is our mechanism:}
\begin{enumerate}
  \item The effect gets \emph{worse} with more data (42\% degradation
    at $N=2000$ vs.\ peak at $N=50$). Regressional Goodhart would get
    \emph{better} with more data.
  \item The direction of degradation is predictable from the cosine
    similarity: diversity degrades because $\cos(\mathbf{w}_d,
    \hat{\mathbf{w}}_u) < 0$.
  \item Strategy~1 (better specification, $\sigma \to 0$) \emph{also}
    harms diversity, confirming the effect is about the inherent
    opposition between stakeholders, not about specification noise.
\end{enumerate}

\subsection{Causal Goodhart}

\textbf{The everyday version.} A hospital measures recovery rates as a
proxy for care quality. Hospital A starts refusing to admit severely ill
patients. Recovery rates soar. Care quality doesn't change---the
hospital just gamed the input distribution.

More precisely: recovery rate \emph{correlates} with care quality in the
observational data. But the causal structure is:
\begin{center}
  care quality $\to$ recovery \\
  patient severity $\to$ recovery
\end{center}
Optimizing recovery by manipulating patient severity breaks the
correlation without changing care quality. The proxy captured an
association, not a causal relationship.

\textbf{The math.} Let $R^*$ be the true reward and $X$ a confounding
variable:
\begin{equation}
  \hat{R}(\pi) = h(R^*(\pi), X(\pi))
\end{equation}
In observational data, $\hat{R}$ and $R^*$ are correlated because $X$
covaries with both. But the optimizer can change $X$ independently of
$R^*$, breaking the correlation.

Formally, the problem is that
$\arg\max_\pi \hat{R}(\pi) \neq \arg\max_\pi R^*(\pi)$
because the optimizer exploits the $X$ pathway rather than the $R^*$
pathway.

\textbf{Key properties:}
\begin{itemize}
  \item Requires an \emph{actionable} confounder---a variable the
    optimizer can manipulate that affects the proxy but not the truth.
  \item Interventional data (randomized experiments) break the
    confounding. Observational proxies are vulnerable; causal proxies
    are not.
  \item This is arguably the most common form of Goodhart in real
    organizations (gaming metrics by manipulating inputs).
\end{itemize}

\textbf{In our setting:} Causal Goodhart is \emph{not} the primary
mechanism. Our BT model operates on genre features, which are fixed
properties of movies---the optimizer can't manipulate them. There's no
confounder being exploited. The model selects content based on fixed
features, not by gaming an input distribution.

However, in a \emph{production} recommendation system, causal Goodhart
is very real: creators gaming the algorithm (e.g., adding trending
keywords to boost visibility without improving content quality). Our
synthetic benchmark doesn't capture this because the content pool is
fixed.

\subsection{Adversarial Goodhart}

\textbf{The everyday version.} A teacher designs a test (proxy for
knowledge). A student obtains the answer key and memorizes it. The
student scores perfectly on the test while learning nothing. The proxy
was designed in good faith, but an adversary deliberately exploited it.

\textbf{The math.} An adversary $\mathcal{A}$ observes the proxy
$\hat{R}$ and constructs a policy $\pi_{\mathcal{A}}$ that maximizes
$\hat{R}$ while minimizing $R^*$:
\begin{equation}
  \pi_{\mathcal{A}} = \arg\max_\pi \left[
    \hat{R}(\pi) - \lambda R^*(\pi)
  \right]
\end{equation}
This is the worst case: the misalignment is intentional and targeted.

\textbf{Key properties:}
\begin{itemize}
  \item Requires an intelligent adversary, not just a misaligned proxy.
  \item The severity is unbounded---the adversary can always find the
    worst-case exploitation.
  \item Defenses require robustness guarantees, not just better proxies.
\end{itemize}

\textbf{In our setting:} Not applicable. Our stakeholders are passive
utility functions, not strategic agents. No one is deliberately trying
to exploit the reward model. (Though in production platforms,
adversarial Goodhart from bot networks and engagement farms is very
real.)

\subsection{Summary: Which Mechanism Drives Our Results?}

\begin{center}
\begin{tabular}{lccl}
\toprule
\textbf{Mechanism} & \textbf{Data helps?} & \textbf{In our setting?} & \textbf{Evidence} \\
\midrule
Regressional & Yes ($\downarrow$ with $N$) & Minor
  & Held-out gap $< 2\%$ \\
\textbf{Extremal} & \textbf{No ($\uparrow$ with $N$)} & \textbf{Dominant}
  & \textbf{42\% degradation, $\cos < 0$} \\
Causal & Depends & Not in benchmark
  & Fixed content pool \\
Adversarial & N/A & No
  & No strategic agents \\
\bottomrule
\end{tabular}
\end{center}

The multi-stakeholder setting makes extremal Goodhart especially sharp:
the ``extreme region'' where the proxy diverges from truth isn't some
obscure corner of policy space. It's the \emph{entire region where
content is user-optimal but diversity-poor}---and BT is designed to
find exactly that region as $N$ increases.

This is what makes our contribution distinct from Skalse et al.: in
their single-objective framework, you need the proxy to be imperfect in
a specific way for extremal Goodhart to occur. In our multi-stakeholder
framework, the ``imperfection'' is \emph{structural}---it's the angle
between stakeholder utilities. As long as $\cos(\mathbf{w}_{\text{target}},
\mathbf{w}_{\text{hidden}}) < 0$, more data \emph{always} hurts the
hidden stakeholder. The Goodhart effect is not a pathology---it's a
geometric consequence of opposing utilities.

\end{document}
